\Chapter{51}

\Verse{1}
{
    \zA{话说众人闻得宝琴将素昔所经过各省内古迹为题，做了十首怀古绝句，内隐十 物，皆说：“这自然新巧。” 都争着看时，只见写道是： 赤壁怀古
        赤壁沉埋水不流，徒留名姓载空舟。 喧阗一炬悲风冷，无限英魂在内游。 交趾怀古 铜柱金城振纪纲，声传海外播戎羌。 马援自是功劳大，铁笛无烦说子房。
        钟山怀古 名利何曾伴女身，无端被诏出凡尘。 牵连大抵难休绝，莫怨他人嘲笑频。 淮阴怀古 壮士须防恶犬欺，三齐位定盖棺时。
        寄言世俗休轻鄙，一饭之恩死也知。}
}
{
    \eA{When the party heard, the story goes, that Pao-ch'in had made the old places
        of interest she had, in days gone by, visited in the various provinces, the
        theme of her verses, and that she had composed ten stanzas with four lines
        in each, which though referring to relics of antiquity, bore covertly on ten
        common objects, they all opined that they must be novel and ingenious, and
        they vied with each other in examining the text.}
}
{
}

\Verse{2}
{
    \zA{广陵怀古 蝉噪鸦栖转眼过，隋堤风景近如何？ 只缘占尽风流号，惹得纷纷口舌多。 桃叶渡怀古 衰草闲花映浅池，桃枝桃叶总分离。
        六朝梁栋多如许，小照空悬壁上题。 青冢怀古 黑水茫茫咽不流，冰弦拨尽曲中愁。 汉家制度诚堪笑，樗栎应惭万古羞。 马嵬怀古
        寂寞脂痕积汗光，温柔一旦付东洋。 只因遗得风流迹，此日衣裳尚有香。 蒲东寺怀古 小红骨贱一身轻，私掖偷携强撮成。 虽被夫人时吊起，已经勾引彼同行。
        梅花观怀古 不在梅边在柳边，个中谁拾画婵娟？}
}
{
    \eA{On perusal, they read: On the relics of Ch'ih Pi:  Deep in Ch'ih Pi doth
        water lie concealed which does not onward flow. There but remains a name and
        surname contained in an empty boat. When with a clamorous din the fire
        breaks out, the sad wind waxes  cold. An endless host of eminent spirits
        wander about inside. On the ancient remains in Chiao Chih:  Posts of copper
        and walls of gold protect the capital.}
}
{
}

\Verse{3}
{
    \zA{团圆莫忆春香到，一别西风又一年。 众人看了，都称奇妙。 宝钗先说道：“前八首都是史鉴上有据的，后二首却无 考。 我们也不大懂得，不如另做两首为是。”
        黛玉忙拦着：“这宝姐姐也忒胶柱鼓 瑟、矫揉造作了。 两首虽于史鉴上无考，咱们虽不曾看这些外传，不知底里，难道 咱们连两本戏也没见过不成？
        那三岁的孩子也知道，何况咱们？” 探春便道：“这 话正是了。” 李纨又道：“况且他原走到这个地方的。}
}
{
    \eA{Its fame is spread beyond the seas, scattered in foreign lands. How true it
        is that Ma Yüan's achievements have been great. The flute of iron need not
        trouble to sing of Tzu Fang. On the vestiges of former times in Chung Shan:
        Renown and gain do they, at any time, fall to a woman's share? For no reason
        have I been bidden come into the mortal world. How hard a task, in point of
        fact, it is to stop solicitude!}
}
{
}

\Verse{4}
{
    \zA{这两件事虽无考，古往今来， 以讹传讹，好事者竟故意的弄出这古迹来以愚人。 比如那年上京的时节，便是关夫 子的坟，倒见了三四处。
        关夫子一身事业皆是有据的，如何又有许多的坟？ 自然是 后来人敬爱他生前为人，只怕从这敬爱上穿凿出来也是有的。 及至看《广舆记》上，
        不止关夫子的坟多有，古来有名望的人，那坟就不少。 无考的古迹更多。 如今这两 首诗虽无考，凡说书唱戏，甚至于求的签上都有。 老少男女俗语口头，人人皆知皆
        说的。}
}
{
    \eA{Don't bear a grudge against such people as may oft times jeer at you! On
        things of historic interest in Huai Yin:  The sturdy man must ever mind the
        insults of the vicious dog. Th' official's rank in San Ch'i was but fixed
        when his coffin was  closed  Tell all people that upon earth do dwell to
        look down upon none. The bounty of one single bowl of rice should be
        treasured till death.}
}
{
}

\Verse{5}
{
    \zA{况且又并不是看了《西厢记》、《牡丹亭》的词曲，怕看了邪书了。 这也无 妨，只管留着。” 宝钗听说，方罢了。 大家猜了一回，皆不是的。
        冬日天短，觉得又是吃晚饭时候，一齐往前头来吃晚饭。 因有人回王夫人说： “袭人的哥哥花自芳，在外头回进来说，他母亲病重了，想他女儿。 他来求恩典，
        接袭人家去走走。” 王夫人听了，便说：“人家母女一场，岂有不许他去的呢。” 一面就叫了凤姐来告诉了，命他酌量办理。}
}
{
    \eA{On events of old in Kuang Lin:  Cicadas chirp; crows roost; but, in a
        twinkle, they are gone. How fares these latter days the scenery in Sui T'i?
        It's all because he has so long enjoyed so fine a fame,  That he has given
        rise around to so many disputes. On the ancient remains of the T'ao Yeh
        ferry:  Dry grass and parchèd plants their reflex cast upon the shallow
        pond.}
}
{
}

\Verse{6}
{
    \zA{凤姐儿答应了，回至屋里，便命周瑞家 的去告诉袭人原故。 吩咐周瑞家的：“再将跟着出门的媳妇传一个，你们两个人， 再带两个小丫头子，跟了袭人去。
        分头派四个有年纪的跟车。 要一辆大车，你们带 着坐，一辆小车，给丫头们坐。” 周瑞家的答应了，才要去，凤姐又道：“那袭人
        是个省事的，你告诉说我的话：叫他穿几件颜色好衣裳，大大的包一包袱衣裳拿着， 包袱要好好的，拿手炉也拿好的。 临走时，叫他先到这里来我瞧。” 周瑞家的答应
        去了。}
}
{
    \eA{The peach tree branches and peach leaves will bid farewell at last. What a
        large number of structures in Liu Ch'ao raise their heads. A small picture
        with a motto hangs on the hollow wall. On the antique vestiges of Ch'ing
        Chung:  The black stream stretches far and wide, but hindered is its course.
        What time were no more thrummed the frozen cords, the songs waxed  sad.}
}
{
}

\Verse{7}
{
    \zA{半日，果见袭人穿戴了，两个丫头和周瑞家的拿着手炉和衣包，凤姐看袭人头 上戴着几枝金钗珠钏，倒也华丽，又看身上穿着桃红百花刻丝银鼠袄，葱绿盘金彩
        绣绵裙，外面穿着青缎灰鼠褂。 凤姐笑道：“这三件衣裳都是太太的，赏了你倒是 好的。 但这褂子太素了些，如今穿着也冷，你该穿一件大毛的。” 袭人笑道：“太
        太就给了这件灰鼠的，还有件银鼠的。 说赶年下再给大毛的呢。” 凤姐笑道：“我 倒有一件大毛的，我嫌风毛出的不好了正要改去，也罢，先给你穿去罢。}
}
{
    \eA{The policy of the Han dynasty was in truth strange! A worthless officer must
        for a thousand years feel shame. On things of historic renown in Ma Wei:
        Quiet the spots of rouge with sweat pile up and shine. Gentleness in a
        moment vanishes and goes. It is because traces remain of his fine looks,
        That to this day his clothes a fragrance still emit.}
}
{
}

\Verse{8}
{
    \zA{等年下太 太给你做的时节，我再改罢。 只当你还我的一样。” 众人都笑道：“奶奶惯会说这 话。
        成年家大手大脚的，替太太不知背地里赔垫了多少东西，真真赔的是说不出来 的，那里又和太太算去？ 偏这会子又说这小气话取笑来了。” 凤姐儿笑道：“太太
        那里想的到这些？ 究竟这又不是正经事。 再不照管，也是大家的体面； 说不得我自 己吃些亏，把众人打扮体统了，宁可我得个好名儿也罢了。
        一个一个‘烧糊了的？}
}
{
    \eA{On events of the past connected with the Pu Tung temple:  The small red lamp
        is wholly made of thin bone, and is light. Furtively was it brought along
        but by force was it stol'n. Oft was it, it is true, hung by the mistress'
        own hands,  But long ere this has she allured it to speed off with her. On
        the scenery about the Mei Hua (Plum Bloom) monastery. If not by the plum
        trees, then by the willows it must be.}
}
{
}

\Verse{9}
{
    \zA{子’似的，人先笑话我，说我当家倒把人弄出个花子来了。” 众人听了，都叹说： “谁似奶奶这么着圣明，在上体贴太太，在下又疼顾下人。” 一面说，一面只见凤
        姐命平儿将昨日那件石青刻丝八团天马皮褂子拿出来，给了袭人。 又看包袱，只得 一个弹墨花绫水红绸里的夹包袱，里面只见包着两件半旧绵袄合皮褂子。 凤姐又命
        平儿把一个玉色绸里的哆罗呢包袱拿出来，又命包上一件雪褂子。 平儿走去拿了出来，一件是件旧大红猩猩毡的，一件是半旧大红羽缎的。}
}
{
    \eA{Has any one picked up in there the likeness of a girl? Don't fret about
        meeting again; in spring its scent returns. Soon as it's gone, and west
        winds blow, another year has flown. When the party had done reading the
        verses, they with perfect unanimity extolled their extraordinary excellence.
        Pao-ch'ai was, however, the first to raise any objections.}
}
{
}

\Verse{10}
{
    \zA{袭人 道：“一件就当不起了。” 平儿笑道：“你拿这猩猩毡的。 把这件顺手带出来，叫
        人给邢大姑娘送去，昨儿那么大雪，人人都穿着不是猩猩毡、就是羽缎的，十来件 大红衣裳，映着大雪，好不齐整。 只有他穿着那几件旧衣裳，越发显的拱肩缩背，
        好不可怜见的，如今把这件给他罢。” 凤姐笑道：“我的东西，他私自就要给人。 我一个还花不够，再添上你提着，更好了！” 众人笑道：“这都是奶奶素日孝敬太
        太，疼爱下人。}
}
{
    \eA{"The first eight stanzas," she said, "are founded upon the testimony of the
        historical works. But as for the last two stanzas, there's no knowing where
        they come from. Besides, we don't quite fathom their meaning. Wouldn't it be
        better then if two other stanzas were written?" Tai-yü hastened to interrupt
        her.}
}
{
}

\Verse{11}
{
    \zA{要是奶奶素日是小气的，收着东西为事的，不顾下人的，姑娘那里 敢这么着？” 凤姐笑道：“所以知道我的，也就是他还知三分罢了。” 说着，又嘱
        咐袭人道：“你妈要好了就罢，要不中用了，只得住下，打发人来回我，我再另打 发人给你送铺盖去。 可别使他们的铺盖和梳头的家伙。” 又吩咐周瑞家的道：“你
        们自然是知道这里的规矩的，也不用我吩咐了。” 周瑞家的答应：“都知道：我们 这去到那里，总叫他们的人回避。 要住下，必是另要一两间内房的。”}
}
{
    \eA{"The lines composed by cousin Pao ch'in are indeed devised in a too
        pigheaded and fast-and-loose sort of way," she observed. "The two stanzas
        are, I admit, not to be traced in the historical works, but though we've
        never read such outside traditions, and haven't any idea what lies at the
        bottom of them, have we not likely seen a couple of plays?}
}
{
}

\Verse{12}
{
    \zA{说着，跟了 袭人出去，又吩咐小厮预备灯笼，遂坐车往花自芳家来，不在话下。 这里凤姐又将怡红院的嬷嬷唤了两个来，吩咐道：“袭人只怕不来家了。 你们
        素日知道那个大丫头知好歹，派出来在宝玉屋里上夜。 你们也好生照管着，别由着 宝玉胡闹。” 两个嬷嬷答应着去了，一时来回说：“派了晴雯和麝月在屋里，我们
        四个人原是轮流着带管上夜的。” 凤姐听了点头，又说道：“晚上催他早睡，早上 催他早起。” 老嬷嬷们答应了，自回园去。}
}
{
    \eA{What child of three years old hasn't some notion about them, and how much
        more such as we?" "What she says is perfectly correct," T'an Ch'un chimed
        in. "She has besides," Li Wan then remarked, "been to these places herself.}
}
{
}

\Verse{13}
{
    \zA{一时果有周瑞家的带了信回凤姐说：“袭 人之母业已停床，不能回来。” 凤姐回明了王夫人，一面着人往大观园去取他的铺 盖妆奁。
        宝玉看着晴雯麝月二人打点妥当。 送去之后，晴雯麝月皆卸罢残妆，脱换过裙袄。 晴雯只在熏笼上围坐，麝月笑 道：“你今儿别装小姐了，我劝你也动一动儿。”
        晴雯道：“等你们都去净了，我 再动不迟。 有你们一日，我且受用一日。” 麝月笑道：“好姐姐，我铺床，你把那 穿衣镜的套子放下来，上头的划子划上。
        你的身量比我高些。”}
}
{
    \eA{But though there be no mention anywhere of these two references, falsehoods
        have from old till now been propagated, and busybodies have, in fact,
        intentionally invented such relics of ancient times with a view of
        bamboozling people. That year, for instance, in which we travelled up here
        to the capital, we came across graves raised to Kuan, the sage, in three or
        four distinct places.}
}
{
}

\Verse{14}
{
    \zA{说着，便去给宝玉 铺床。 晴雯？ 了一声，笑道：“人家才坐暖和了，你就来闹。” 此时宝玉正坐着纳
        闷，想袭人之母不知是死是活，忽听见晴雯如此说，便自己起身出去，放下镜套， 划上消息。 进来笑道：“你们暖和罢，我都弄完了。” 晴雯笑道：“终久暖和不成，
        我又想起来，汤婆子还没拿来呢。” 麝月道：“这难为你想着!他素日又不要汤壶， 咱们那熏笼上又暖和，比不得那屋里炕凉，今儿可以不用。”}
}
{
    \eA{Now the circumstances of the whole existence of Kuan the sage are
        established by actual proof, so how could there again in his case exist a
        lot of graves? This must arise from the esteem in which he is held by
        posterity for the way he acquitted himself of his duties during his
        lifetime. And it is presumably to this esteem that this fiction owes its
        origin. This is quite possible enough.}
}
{
}

\Verse{15}
{
    \zA{宝玉笑道：“你们两 个都在那上头睡了，我这外边没个人，我怪怕的，一夜也睡不着。” 晴雯道：“我 是在这里睡的，麝月，你叫他往外边睡去。”
        说话之间，天已一更，麝月早已放下 帘幔，移灯炷香，伏侍宝玉卧下，二人方睡。 晴雯自在熏笼上，麝月便在暖阁外边。 至三更以后，宝玉睡梦之中，便叫袭人。
        叫了两声，无人答应，自己醒了，方 想起袭人不在家，自己也好笑起来。 晴雯已醒，因唤麝月道：“连我都醒了，他守 在旁边还不知道，真是挺死尸呢！”}
}
{
    \eA{Even in the 'Kuang Yü Chi', you will see that not only are numerous tombs of
        the sage Kuan spoken of, but that bygone persons of note are assigned tombs
        not few in number. But there are many more relics of antiquity, about which
        no testimony can be gathered. The matter treated in the two stanzas, now in
        point, is, of course, not borne out by any actual record;}
}
{
}

\Verse{16}
{
    \zA{麝月翻身打个哈什，笑道：“他叫袭人，与我 什么相干！” 因问：“做什么？” 宝玉说要吃茶。 麝月忙起来，单穿着红绸小绵袄 儿。
        宝玉道：“披了我的皮袄再去，仔细冷着。” 麝月听说，回手便把宝玉披着起 来的一件貂颏满襟暖袄披上，下去向盆内洗洗手，先倒了一钟温水，拿了大漱盂，
        宝玉漱了口。 然后才向茶桶上取了茶碗，先用温水过了，向暖壶中倒了半碗茶，递 给宝玉吃了，自己也漱了一漱，吃了半碗。
        晴雯笑道：“好妹妹，也赏我一口儿呢。”}
}
{
    \eA{yet in every story, that is told, in every play, that is sung, and on the
        various slips as well used for fortune telling, it is invariably to be
        found. Old and young, men and women, do all understand it and speak of it,
        whether in proverbs or in their everyday talk.}
}
{
}

\Verse{17}
{
    \zA{麝月笑道：“越发上脸儿了！” 晴雯道：“好妹妹，明儿晚上你别动，我伏侍你一 夜，如何？” 麝月听说，只得也伏侍他漱了口，倒了半碗茶给他吃了。 麝月笑道：
        “你们两个别睡，说着话儿，我出去走走回来。” 晴雯笑道：“外头有个鬼等着呢。” 宝玉道：“外头自然有大月亮的。 我们说着话，你只管去。”
        一面说，一面便嗽了 两声。 麝月便开了后房门，揭起毡帘一看，果然好月色。}
}
{
    \eA{They don't resemble, besides, the ballads encountered in the 'Hsi Hsiang
        Chi,' and 'Mou Tan T'ing,' to justify us to fear that we might be setting
        eyes upon some corrupt text. They are quite harmless; so we'd better keep
        them!" Pao-ch'ai, after these arguments, dropped at length all discussion.
        They thereupon tried for a time to guess the stanzas. None, however, of
        their solutions turned out to be correct.}
}
{
}

\Verse{18}
{
    \zA{晴雯等他出去，便欲唬他 玩耍，仗着素日比别人气壮，不畏寒冷，也不披衣，只穿着小袄便蹑手蹑脚的下了 熏笼，随后出来。 宝玉劝道：“罢呀，冻着不是玩的！”
        晴雯只摆手，随后出了屋 门，只见月光如水。 忽听一阵微风，只觉侵肌透骨，不禁毛骨悚然。 心下自思道： “怪道人说热身子不可被风吹，这一冷果然利害。”
        一面正要唬他，只听宝玉在内 高声说道：“晴雯出来了！” 晴雯忙回身进来，笑道：“那里就唬死了他了？ 偏惯会这么蝎蝎螫螫老婆子的 样儿。”}
}
{
    \eA{But as the days in winter are short, and they saw that it was time for their
        evening meal, they adjourned to the front part of the compound for their
        supper.}
}
{
}

\Verse{19}
{
    \zA{宝玉笑道：“倒不是怕唬坏了他。 头一件你冻着也不好，二则他不防，不 免一喊，倘或惊醒了别人，不说咱们是玩意儿，倒反说袭人才去了一夜，你们就见 神见鬼的。
        你来把我这边的被掖掖罢。” 晴雯听说，就上来掖了一掖，伸手进去就 渥一渥。 宝玉笑道：“好冷手，我说看冻着。” 一面又见晴雯两腮如胭脂一般，用
        手摸一摸，也觉冰冷。 宝玉道：“快进被来渥渥罢。”}
}
{
    \eA{The servants at this stage announced to Madame Wang that Hsi Jen's elder
        brother, Hua Tzu-fang, was outside, and reported to her that he had entered
        the city to say that his mother was lying in bed dangerously ill, and that
        she was so longing to see her daughter that he had come to beg for the
        favour of taking Hsi Jen home on a visit.}
}
{
}

\Verse{20}
{
    \zA{一语未了，只听咯噔的一声 门响，麝月慌慌张张的笑着进来，说着笑道：“唬我一跳好的!黑影子里，山子石 后头，只见一个人蹲着。
        我才要叫喊，原来是那个大锦鸡，见了人，一飞飞到亮处 来，我才见了。 要冒冒失失一嚷，倒闹起人来。” 一面说，一面洗手，又笑道：“说 晴雯出去了？
        我怎么没见。 一定是要唬我去了。” 宝玉笑道：“这不是他？ 在这里渥 着呢。 我若不嚷的快，可是倒唬一跳。” 晴雯笑道：“也不用我唬去，这小蹄子已
        经自惊自怪的了。”}
}
{
    \eA{As soon as Madame Wang heard the news, she dilated for a while upon people's
        mothers and daughters, and of course she did not withhold her consent.
        Sending therefore at the same time for lady Feng, she communicated the
        tidings to her, and enjoined her to deliberate, and take suitable action.}
}
{
}

\Verse{21}
{
    \zA{一面说，一面仍回自己被中去。 麝月道：“你就这么‘跑解马’ 的打扮儿，伶伶俐俐的出去了不成？” 宝玉笑道：“可不就是这么出去了。” 麝月
        道：“你死不拣好日子!你出去自站一站，瞧把皮不冻破了你的。” 说着又将火盆 上的铜罩揭起，拿灰锹重将熟炭埋了一埋，拈了两块速香放上，仍旧罩了。 至屏后，
        重剔亮了灯，方才睡下。 晴雯因方才一冷，如今又一暖，不觉打了两个嚏喷。 宝玉叹道：“如何？ 到底 伤了风了。”}
}
{
    \eA{Lady Feng signified her willingness to do what was necessary, and, returning
        to her quarters, she there and then commissioned Chou Jui's wife to go and
        break the news to Hsi Jen. "Send also," she went on to direct Mrs. Chou,
        "for one of the married-women, who are in attendance when we go
        out-of-doors, and let you two, together with a couple of young maids, follow
        Hsi Jen home.}
}
{
}

\Verse{22}
{
    \zA{麝月笑道：“他早起就嚷不受用，一日也没吃碗正经饭。 他这会子不 说保养着些，还要捉弄人，明儿病了，叫他自作自受。” 宝玉问道：“头上热不热？”
        晴雯嗽了两声，说道：“不相干，那里这么娇嫩起来了。” 说着，只听外间屋里？ 上的自鸣钟“当当”的两声，外间值宿的老嬷嬷嗽了两声，因说道：“姑娘们睡罢，
        明儿再说笑罢。” 宝玉方悄悄的笑道：“咱们别说话了，看又惹他们说话。” 说着， 方大家睡了。 至次日起来，晴雯果觉有些鼻塞声重，懒怠动弹。}
}
{
    \eA{But despatch four cart attendants, well up in years, to look everywhere for
        a spacious curricle for you as well as her, and a small carriage for the
        maids." "All right!" acquiesced Chou Jui's wife. But just as she was about
        to start, lady Feng continued her injunctions. "Hsi Jen," she added; "is a
        person not fond of any fuss, so tell her that it's I who have given the
        orders;}
}
{
}

\Verse{23}
{
    \zA{宝玉道：“快别声张。 太太 知道了，又要叫你搬回家去养着。 家里纵好，到底冷些，不如在这里。 你就在里间
        屋里躺着，我叫人请了大夫，悄悄的从后门进来瞧瞧就是了。” 晴雯道：“虽这么 说，你到底要告诉大奶奶一声儿。 不然一时大夫来了，人问起来怎么说呢？” 宝玉
        听了有理，便唤一个老嬷嬷来吩咐道：“你回大奶奶去，就说晴雯白冷着了些，不 是什么大病。 袭人又不在家，他若家去养病，这里更没有人了。}
}
{
    \eA{and impress upon her that she must put on several nice, coloured clothes,
        and pack up a large valise full of wearing apparel. Her valise, must be a
        handsome one; and she must take a decent hand-stove. Bid her too first come
        and look me up here when she's about to start." Mrs. Chou promised to
        execute her directions and went on her way.}
}
{
}

\Verse{24}
{
    \zA{传一个大夫，从后 门悄悄的进来瞧瞧，别回太太了。” 老嬷嬷去了，半日回来说：“大奶奶知道了。 说两剂药好了便罢，若不好时，还是出去为是。
        如今的时气不好，沾染了别人事小， 姑娘们的身子要紧。” 晴雯睡在暖阁里，只管咳嗽，听了这话，气的嚷道：“我那 里就害瘟病了？ 生怕招了人。
        我离了这里，看你们这一辈子都别头疼脑热的！” 说 着，便真要起来。 宝玉忙按他，笑道：“别生气，这原是他的责任，生恐太太知道 了说他。 不过白说一句。}
}
{
    \eA{After a long interval, (lady Feng) actually saw Hsi Jen arrive, got up in
        full costume and head-gear, and with her two waiting-maids and Chou Jui's
        wife, who carried the hand-stove and the valise packed up with clothes. Lady
        Feng's eye was attracted by several golden hairpins and pearl ornaments of
        great brilliancy and beauty, which Hsi Jen wore in her coiffure.}
}
{
}

\Verse{25}
{
    \zA{你素昔又爱生气，如今肝火自然又盛了。” 正说时，人回大夫来了。 宝玉便走过来，避在书架后面。 只见两三个后门口的 老婆子带了一个太医进来。
        这里的丫头都回避了，有三四个老嬷嬷放下暖阁上的大 红绣幔，晴雯从幔中单伸出手来。 那大夫见这只手上有两根指甲，足有二三寸长，
        尚有金凤仙花染的通红的痕迹，便回过头来。 有一个老嬷嬷忙拿了一块绢子掩上 了。 那大夫方诊了一回脉，起身到外间，向嬷嬷们说道：“小姐的症是外感内滞。}
}
{
    \eA{Her gaze was further struck by the peach-red stiff silk jacket she had on,
        brocaded with all sorts of flowers and lined with ermine, by her leek-green
        wadded jupe, artistically ornamented with coils of gold thread, and by the
        bluish satin and grey squirrel pelisse she was wrapped in. "These three
        articles of clothing, given to you by our dowager lady," lady Feng smiled,
        "are all very nice;}
}
{
}

\Verse{26}
{
    \zA{近日时气不好，竟算是个小伤寒。 幸亏是小姐，素日饮食有限，风寒也不大，不过 是气血原弱，偶然沾染了些，吃两剂药疏散疏散就好了。” 说着，便又随婆子们出
        去。 彼时李纨已遣人知会过后门上的人及各处丫鬟回避。 大夫只见了园中景致，并 不曾见一个女子。 一时出了园门，就在守园门的小厮们的班房内坐了，开了药方。
        老嬷嬷道：“老爷且别去，我们小爷罗嗦，恐怕还有话问。” 那太医忙道：“方才 不是小姐，是位爷不成？}
}
{
    \eA{but this pelisse is somewhat too plain. If you wear this, you'll besides
        feel cold, so put on one with long fur." "Our Madame Wang," Hsi Jen
        laughingly rejoined, "gave me this one with the grey squirrel. I've also got
        one with ermine. She says that when the end of the year draws nigh, she'll
        let me have one with long fur." "I've got one with long fur," lady Feng
        proceeded with a smile.}
}
{
}

\Verse{27}
{
    \zA{那屋子竟是绣房，又是放下幔子来瞧的，如何是位爷呢？” 老嬷嬷笑道：“我的老爷，怪道小子才说：‘今儿请了一位新太医来了。’ 真不知 我们家的事。
        那屋子是我们小哥儿的，那人是屋里的丫头，倒是个‘大姐’，那里 的小姐的绣房？ 小姐病了，你那么容易就进去了？” 说着，拿了药方进去。
        宝玉看时，上面有紫苏、桔梗、防风、荆芥等药，后面又有枳实、麻黄。 宝玉 道：“该死该死，他拿着女孩儿们也像我们一样的治法，如何使得？}
}
{
    \eA{"I don't fancy it much as the fringe does not hang with grace. I was on the
        point of having it changed; but, never mind, I'll let you first use it; and,
        when at the close of the year, Madame Wang has one made for you, I can then
        have mine altered, and it will come to the same thing as if you were
        returning it like that to me." One and all laughed. "That's the way of
        talking into which her ladyship has got!"}
}
{
}

\Verse{28}
{
    \zA{凭他有什么内 滞，这枳实、麻黄如何禁得？ 谁请了来的？ 快打发他去罢，再请一个熟的来罢。” 老 嬷嬷道：“用药好不好，我们不知道。
        如今再叫小厮去请王大夫去倒容易，只是这 个大夫又不是告诉总管房请的，这马钱是要给他的。” 宝玉道：“给他多少？” 婆
        子道：“少了不好，看来得一两银子，才是我们这样门户的礼。” 宝玉道：“王大 夫来了，给他多少？”}
}
{
    \eA{they observed. "There she is the whole year round recklessly carelessly and
        secretly making good, on Madame Wang's account, ever so many things; how
        many there is no saying; for really the things for which compensation is
        made, cannot be so much as enumerated; and does she ever go, and settle
        scores with Madame Wang?}
}
{
}

\Verse{29}
{
    \zA{婆子笑道：“王大夫和张大夫每常来了，也并没个给钱的， 不过每年四节一个趸儿送礼，那是一定的年例。 这个人新来了一次，须得给他一两 银子。”
        宝玉听说，就命麝月去取银子。 麝月道：“花大姐姐还不知搁在那里呢？” 宝玉道：“我常见着在那小螺甸柜子里拿银子，我和你找去。” 说着二人来至袭人
        堆东西的屋内，开了螺甸柜子。 上一？ 都是些笔墨、扇子、香饼、各色荷包、汗 巾等类的东西，下一？ 却有几串钱。}
}
{
    \eA{and here she comes, on this occasion, and gives vent again to this mean
        language, in order to poke fun at people!" "How could Madame Wang," lady
        Feng laughed, "ever give a thought to such trifles as these? They are, in
        fact, matters of no consequence. Yet were I not to look after them, it would
        be a disgrace to all of us, and needless to say, I would myself get into
        some scrape.}
}
{
}

\Verse{30}
{
    \zA{于是开了抽屉，才看见一个小笸箩内放着几块 银子，倒也有戥子。 麝月便拿了一块银，提起戥子来问宝玉：“那是一两的星儿？”
        宝玉笑道：“你问的我有趣儿，你倒成了是才来的了。” 麝月也笑了，又要去问人。 宝玉道：“拣那大的给他一块就是了。 又不做买卖，算这些做什么。”
        麝月听了， 便放下戥子，拣了一块掂了一掂，笑道：“这一块只怕是一两了。 宁可多些好，别 少了叫那穷小子笑话：不说咱们不认得戥子，倒说咱们有心小气似的。”}
}
{
    \eA{It's far better that I should dress you all properly, and so get a fair name
        and finish; for were each of you to cut the figure of a burnt cake, people
        would first and foremost ridicule me, by saying that in looking after the
        household I have, instead of doing good, been the means of making beggars of
        you!" After hearing her out, the whole party heaved a sigh.}
}
{
}

\Verse{31}
{
    \zA{那婆子站 在门口笑道：“那是五两的锭子夹了半个，这一块至少还有二两呢。 这会子又没夹 剪，姑娘收了这块，拣一块小些的。”
        麝月早关了柜子出来，笑道：“谁又找去呢， 多少你拿了去就完了！” 宝玉道：“你快叫焙茗再请个大夫来罢。” 婆子接了银子， 自去料理。
        一时焙茗果请了王大夫来，先诊了脉，后说病症，也与前头不同。 方子上果然 没有枳实、麻黄等药，倒有当归、陈皮、白芍等药，那分两较先也减了些。 宝玉喜
        道：“这才是女孩儿们的药。}
}
{
    \eA{"Who could ever be," they exclaimed, "so intuitively wise as you, to show,
        above, such regard for Madame Wang, and below, such consideration for her
        subordinates?" In the course of these remarks, they noticed lady Feng bid
        P'ing Erh find the dark green stiff silk cloak with white fox, she had worn
        the day before, and give it to Hsi Jen.}
}
{
}

\Verse{32}
{
    \zA{虽疏散，也不可太过。 旧年我病了，却是伤寒，内里 饮食停滞，他瞧了还说我禁不起麻黄、石膏、枳实等狼虎药。 我和你们就如秋天芸
        儿进我的那才开的白海棠似的； 我禁不起的药，你们那里经得起？ 比如人家坟里的 大杨树，看着枝叶茂盛，都是空心子的。” 麝月笑道：“野坟里只有杨树，难道就
        没有松柏不成？ 最讨人嫌的是杨树，那么大树只一点子叶子，没一点风儿他也是乱 响。 你偏要比他，你也太下流了。” 宝玉笑道：“松柏不敢比。}
}
{
    \eA{But perceiving, also, that in the way of a valise, she only had a double one
        made of black spotted, figured sarcenet, with a lining of light red pongee
        silk, and that its contents consisted merely of two wadded jackets, the
        worse for wear, and a pelisse, lady Feng went on to tell P'ing Erh to fetch
        a woollen wrapper, lined with jade-green pongee. But she ordered her besides
        to pack up a snow-cloak for her.}
}
{
}

\Verse{33}
{
    \zA{连孔夫子都说‘岁 寒然后知松柏之后雕’呢，可知这两件东西高雅。 不害臊的才拿他混比呢。” 说着，只见老婆子取了药来。
        宝玉命把煎药的银铞子找了出来，就命在火盆上 煎。 晴雯因说：“正经给他们茶房里煎去罢咧，弄的这屋里药气，如何使得？” 宝
        玉道：“药气比一切的花香还香呢。 神仙采药烧药，再者高人逸士采药治药，最妙 的一件东西。 这屋里我正想各色都齐了，就只少药香，如今恰全了。” 一面说，一
        面早命人煨上。}
}
{
    \eA{P'ing Erh walked away and produced the articles. The one was made of
        deep-red felt, and was old. The other was of deep-red soft satin, neither
        old nor new. "I don't deserve so much as a single one of these," Hsi Jen
        said.}
}
{
}

\Verse{34}
{
    \zA{又嘱咐麝月打点些东西，叫个老嬷嬷去看袭人，劝他少哭。 一一妥 当，方过前边来贾母王夫人处请安吃饭。
        正值凤姐儿和贾母王夫人商议说：“天又短，又冷，不如以后大嫂子带着姑娘 们在园子里吃饭。 等天暖和了，再来回的跑，也不妨。” 王夫人笑道：“这也是好
        主意。 刮风下雪倒便宜。 吃东西受了冷气也不好，空心走来，一肚子冷气，压上些 东西也不好。 不如园子后门里头的五间大屋子，横竖有女人们上夜的，挑两个女厨
        子在那里单给他姐妹弄饭。}
}
{
    \eA{"Keep this felt one for yourself," P'ing Erh smiled, "and take this one
        along with you and tell some one to send it to that elderly girl, who while
        every one, in that heavy fall of snow yesterday, was rolled up in soft
        satin, if not in felt, and while about ten dark red dresses were reflected
        in the deep snow and presented such a fine sight, was the only one attired
        in those shabby old clothes.}
}
{
}

\Verse{35}
{
    \zA{新鲜菜蔬是有分例的，在总管账房里支了去，或要钱要 东西。 那些野鸡獐狍各样野味，分些给他们就是了。” 贾母道：“我也正想着呢， 就怕又添厨房事多些。”
        凤姐道：“并不事多：一样的分例，这里添了，那里减了。 就便多费些事，小姑娘们受了冷气，别人还可，第一，林妹妹如何禁得住？ 就连宝 玉兄弟也禁不住。
        况兼众位姑娘都不是结实身子。” 凤姐儿说毕，未知贾母何言，且听下回分解。 (本章完)}
}
{
    \eA{She seems more than ever to raise her shoulders and double her back. She is
        really to be pitied; so take this now and give it to her!" "She
        surreptitiously wishes to give my things away!" lady Feng laughed. "I
        haven't got enough to spend upon myself and here I have you, better still,
        to instigate me to be more open-handed!"}
}
{
}

\Verse{36}
{
    \zA{}
}
{
    \eA{"This comes from the filial piety your ladyship has ever displayed towards
        Madame Wang," every one laughingly remarked, "and the fond love for those
        below you. For had you been mean and only thought of making much of things
        and not cared a rap for your subordinates, would that girl have presumed to
        behave in this manner?" "If any one therefore has read my heart, it's she,"
        lady Feng rejoined with a laugh, "but yet she only knows it in part." At the
        close of this rejoinder, she again spoke to Hsi Jen. "If your mother gets
        well, all right," she said; "but if anything happens to her, just stay over,
        and send some one to let me know so that I may specially despatch a servant
        to bring you your bedding.}
}
{
}

\Verse{37}
{
    \zA{}
}
{
    \eA{But whatever you do, don't, use their bedding, nor any of their things to
        comb your hair with. As for you people," continuing, she observed to Mrs.
        Chou Jui, "you no doubt are aware of the customs, prevailing in this
        establishment, so that I can dispense with giving you any injunctions."
        "Yes, we know them all," Mrs. Chou Jui assented. "As soon as we get there,
        we'll, of course, request their male inmates to retire out of the way. And
        in the event of our having to stay over, we'll naturally apply for one or
        two extra inner rooms." With these words still on her lips, she followed Hsi
        Jen out of the apartment.}
}
{
}

\Verse{38}
{
    \zA{}
}
{
    \eA{Then directing the servant-boys to prepare the lanterns, they, in due
        course, got into their curricle, and came to Hua Tzu-fang's quarters, where
        we will leave them without any further comment. Lady Feng, meanwhile, sent
        also for two nurses from the I Hung court. "I am afraid," she said to them,
        "that Hsi Jen won't come back, so if there be any elderly girl, who has to
        your knowledge, so far, had her wits about her, depute her to come and keep
        night watch in Pao-yü's rooms. But you nurses must likewise take care and
        exercise some control, for you mustn't let Pao-yü recklessly kick up any
        trouble!"}
}
{
}

\Verse{39}
{
    \zA{}
}
{
    \eA{"Quite so," answered the two nurses, agreeing to her directions, after
        which, they quitted her presence. But not a long interval expired before
        they came to report the result of their search. "We've set our choice upon
        Ch'ing Wen and She Yüeh to put up in his rooms," they reported. "We four
        will take our turn and look after things during the night." When lady Feng
        heard these arrangements, she nodded her head. "At night," she observed,
        "urge him to retire to bed soon; and in the morning press him to get up at
        an early hour." The nurses replied that they would readily carry out her
        orders and returned alone into the garden. In a little time Chou Jui's wife
        actually brought the news, which she imparted to lady Feng, that: "as her
        mother was already beyond hope, Hsi Jen could not come back."}
}
{
}

\Verse{40}
{
    \zA{}
}
{
    \eA{Lady Feng then explained things to Madame Wang, and sent, at the same time,
        servants to the garden of Broad Vista to fetch (Hsi Jen's) bedding and
        toilet effects. Pao-yü watched Ch'ing Wen and She Yüeh get all her
        belongings in proper order. After the things had been despatched, Ch'ing Wen
        and She Yüeh divested themselves of their remaining fineries and changed
        their jupes and jackets. Ch'ing Wen seated herself round a warming-frame.
        "Now," She Yüeh smiled, "you're not to put on the airs of a young lady! I
        advise you to also move about a bit." "When you're all clean gone," Ch'ing
        Wen returned for answer, "I shall have ample time to budge. But every day
        that you people are here, I shall try and enjoy peace and quiet."}
}
{
}

\Verse{41}
{
    \zA{}
}
{
    \eA{"My dear girl," She Yüeh laughed, "I'll make the bed, but drop the cover
        over that cheval-glass and put the catches right; you are so much taller
        than I." So saying, she at once set to work to arrange the bed for Pao-yü.
        "Hai!" ejaculated Ch'ing Wen smiling, "one just sits down to warm one's
        self, and here you come and disturb one!" Pao-yü had at this time been
        sitting, plunged in a despondent mood. The thought of Hsi Jen's mother had
        crossed through his mind and he was wondering whether she could be dead or
        alive, when unexpectedly overhearing Ch'ing Wen pass the remarks she did, he
        speedily sprung up, and came out himself and dropped the cover of the glass,
        and fastened the contrivance, after which he walked into the room. "Warm
        yourselves," he smiled, "I've done all there was to be done."}
}
{
}

\Verse{42}
{
    \zA{}
}
{
    \eA{"I can't manage," Ch'ing Wen rejoined smiling, "to get warm at all. It just
        also strikes me that the warming-pan hasn't yet been brought." "You've had
        the trouble to think of it!" She Yüeh observed. "But you've never wanted a
        chafing-dish before. It's so warm besides on that warming-frame of ours; not
        like the stove-couch in that room, which is so cold; so we can very well do
        without it to-day." "If both of you are to sleep on that," Pao-yü smiled,
        "there won't be a soul with me outside, and I shall be in an awful funk.
        Even you won't be able to have a wink of sleep during the whole night!" "As
        far as I'm concerned," Ch'ing Wen put in, "I'm going to sleep in here.
        There's She Yüeh, so you'd better induce her to come and sleep outside."}
}
{
}

\Verse{43}
{
    \zA{}
}
{
    \eA{But while they kept up this conversation, the first watch drew near, and She
        Yüeh at once lowered the mosquito-curtain, removed the lamp, burnt the
        joss-sticks, and waited upon Pao-yü until he got into bed. The two maids
        then retired to rest. Ch'ing Wen reclined all alone on the warming-frame,
        while She Yüeh lay down outside the winter apartments. The third watch had
        come and gone, when Pao-yü, in the midst of a dream, started calling Hsi
        Jen. He uttered her name twice, but no one was about to answer him. And it
        was after he had stirred himself out of sleep that he eventually recalled to
        mind that Hsi Jen was not at home, and he had a hearty fit laughter to
        himself. Ch'ing Wen however had been roused out of her sleep, and she called
        She Yüeh.}
}
{
}

\Verse{44}
{
    \zA{}
}
{
    \eA{"Even I," she said, "have been disturbed, fast asleep though I was; and, lo,
        she keeps a look-out by his very side and doesn't as yet know anything about
        his cries! In very deed she is like a stiff corpse!" She Yüeh twisted
        herself round and yawned. "He calls Hsi Jen," she smilingly rejoined, "so
        what's that to do with me? What do you want?" proceeding, she then inquired
        of him. "I want some tea," Pao-yü replied. She Yüeh hastily jumped out of
        bed, with nothing on but a short wadded coat of red silk. "Throw my pelisse
        over you;" Pao-yü cried; "for mind it's cold!" She Yüeh at these words put
        back her hands, and, taking the warm pelisse, lined even up to the lapel,
        with fur from the neck of the sable, which Pao-yü had put on on getting up,
        she threw it over her shoulders and went below and washed her hands in the
        basin.}
}
{
}

\Verse{45}
{
    \zA{}
}
{
    \eA{Then filling first a cup with tepid water, she brought a large cuspidor for
        Pao-yü to wash his mouth. Afterwards, she drew near the tea-case, and
        getting a cup, she first rinsed it with lukewarm water, and pouring half a
        cup of tea from the warm teapot, she handed it to Pao-yü. After he had done,
        she herself rinsed her mouth, and swallowed half a cupful of tea. "My dear
        girl," Ch'ing Wen interposed smiling, "do give me also a sip." "You put on
        more airs than ever," She Yüeh laughed. "My dear girl;" Ch'ing Wen added,
        "to-morrow night, you needn't budge; I'll wait on you the whole night long.
        What do you say to that?" Hearing this, She Yüeh had no help but to attend
        to her as well, while she washed her mouth, and to pour a cup of tea and
        give it to her to drink.}
}
{
}

\Verse{46}
{
    \zA{}
}
{
    \eA{"Won't you two go to sleep," She Yüeh laughed, "but keep on chatting? I'll
        go out for a time; I'll be back soon." "Are there any evil spirits waiting
        for you outside?" Ch'ing Wen smiled. "It's sure to be bright moonlight out
        of doors," Pao-yü observed, "so go, while we continue our chat." So
        speaking, he coughed twice. She Yüeh opened the back-door, and raising the
        woollen portière and looking out, she saw what a beautiful moonlight there
        really was. Ch'ing Wen allowed her just time enough to leave the room, when
        she felt a wish to frighten her for the sake of fun.}
}
{
}

\Verse{47}
{
    \zA{}
}
{
    \eA{But such reliance did she have in her physique, which had so far proved
        better than that of others, that little worrying her mind about the cold,
        she did not even throw a cloak over her, but putting on a short jacket, she
        descended, with gentle tread and light step, from the warming-frame and was
        making her way out to follow in her wake, when "Hallo!" cried Pao-yü warning
        her. "It's freezing; it's no joke!" Ch'ing Wen merely responded with a wave
        of the hand and sallied out of the door to go in pursuit of her companion.
        The brilliancy of the moon, which met her eye, was as limpid as water. But
        suddenly came a slight gust of wind. She felt it penetrate her very flesh
        and bore through her bones. So much so, that she could not help shuddering
        all over.}
}
{
}

\Verse{48}
{
    \zA{}
}
{
    \eA{"Little wonder is it," she argued within herself, "if people say 'that one
        mustn't, when one's body is warm, expose one's self to the wind.' This cold
        is really dreadful!" She was at the same time just on the point of giving
        (She Yüeh) a start, when she heard Pao-yü shout from inside, "Ch'ing Wen has
        come out." Ch'ing Wen promptly turned back and entered the room. "How could
        I ever frighten her to death?" she laughed. "It's just your way; you're as
        great a coward as an old woman!" "It isn't at all that you might do her harm
        by frightening her," Pao-yü smiled, "but, in the first place, it wouldn't be
        good for you to get frost-bitten; and, in the second, you would take her so
        much off her guard that she won't be able to prevent herself from uttering a
        shout.}
}
{
}

\Verse{49}
{
    \zA{}
}
{
    \eA{So, in the event of rousing any of the others out of their sleep, they won't
        say that we are up to jokes, but maintain instead that just as Hsi Jen is
        gone, you two behave as if you'd come across ghosts or seen evil spirits.
        Come and tuck in the coverlets on this side!" When Ch'ing Wen heard what he
        wanted done she came accordingly and tucked in the covers, and, putting out
        her hands, she inserted them under them, and set to work to warm the
        bedding. "How cold your hand is!" Pao-yü laughingly exclaimed. "I told you
        to look out or you'd freeze!" Noticing at the same time that Ch'ing Wen's
        cheeks were as red as rouge, he rubbed them with his hands.}
}
{
}

\Verse{50}
{
    \zA{}
}
{
    \eA{But as they felt icy cold to his touch, "Come at once under the cover and
        warm yourself!" Pao-yü urged. Hardly, however, had he concluded these words,
        than a sound of 'lo teng' reached their ears from the door, and She Yüeh
        rushed in all in a tremor, laughing the while. "I've had such a fright," she
        smiled, as she went on speaking. "Goodness me! I saw in the black shade, at
        the back of the boulders on that hill, some one squatting, and was about to
        scream, when it turned out to be nothing else than that big golden pheasant.
        As soon as it caught sight of a human being, it flew away. But it was only
        when it reached a moonlit place that I at last found out what it was.}
}
{
}

\Verse{51}
{
    \zA{}
}
{
    \eA{Had I been so heedless as to scream, I would have been the means of getting
        people out of their beds!" Recounting her experiences, she washed her hands.
        "Ch'ing Wen, you say, has gone out," she proceeded laughing, "but how is it
        I never caught a glimpse of her? She must certainly have gone to frighten
        me!" "Isn't this she?" Pao-yü inquired with a smile. "Is she not here
        warming herself? Had I not been quick in shouting, she would verily have
        given you a fright." "There was no need for me to go and frighten her,"
        Ch'ing Wen laughingly observed. "This hussy has frightened her own self."
        With these words she ensconced herself again under her own coverlet. "Did
        you forsooth go out," She Yüeh remarked, "in this smart dress of a
        circus-performer?" "Why, of course, she went out like this!" Pao-yü smiled.}
}
{
}

\Verse{52}
{
    \zA{}
}
{
    \eA{"You wouldn't know, for the life of you, how to choose a felicitous day!"
        She Yüeh added. "There you go and stand about on a fruitless errand. Won't
        your skin get chapped from the frost?" Saying this, she again raised the
        copper cover from the brasier, and, picking up the shovel, she buried the
        live charcoal deep with ashes, and taking two bits of incense of Cambodia
        fragrant wood, she threw them over them. She then re-covered the brasier,
        and repairing to the back of the screen, she gave the lamp a thorough
        trimming to make it throw out more light; after which, she once more laid
        herself down. As Ch'ing Wen had some time before felt cold, and now began to
        get warm again, she unexpectedly sneezed a couple of times. "How about
        that?" sighed Pao-yü. "There you are; you've after all caught a chill!"}
}
{
}

\Verse{53}
{
    \zA{}
}
{
    \eA{"Early this morning," She Yüeh smiled, "she shouted that she wasn't feeling
        quite herself. Neither did she have the whole day a proper bowl of food. And
        now, not to speak of her taking so little care of herself, she is still bent
        upon playing larks upon people! But if she falls ill by and bye, we'll let
        her suffer what she will have brought upon herself." "Is your head hot?"
        Pao-yü asked. "It's nothing at all!" Ch'ing Wen rejoined, after coughing
        twice. "When did I get so delicate?" But while she spoke, they heard the
        striking clock, suspended on the partition wall in the outer rooms, give two
        sounds of 'tang, tang,' and the matron, on the night watch outside, say:
        "Now, young girls, go to sleep. To-morrow will be time enough for you to
        chat and laugh!" "Don't let's talk!"}
}
{
}

\Verse{54}
{
    \zA{}
}
{
    \eA{Pao-yü then whispered, "for, mind, we'll also induce them to start
        chattering." After this, they at last went to sleep. The next day, they got
        up at an early hour. Ch'ing Wen's nose was indeed considerably stopped. Her
        voice was hoarse; and she felt no inclination to move. "Be quick," urged
        Pao-yü, "and don't make a fuss, for your mistress, my mother, may come to
        know of it, and bid you also shift to your house and nurse yourself. Your
        home might, of course, be all very nice, but it's in fact somewhat cold. So
        isn't it better here? Go and lie down in the inner rooms, and I'll give
        orders to some one to send for the doctor to come quietly by the back door
        and have a look at you.}
}
{
}

\Verse{55}
{
    \zA{}
}
{
    \eA{You'll then get all right again." "In spite of what you say," Ch'ing Wen
        demurred, "you must really say something about it to our senior lady, Mrs.
        Chia Chu; otherwise the doctor will be coming unawares, and people will
        begin to ask questions; and what answer could one give them?" Pao-yü found
        what she said so full of reason that he called an old nurse. "Go and deliver
        this message to your senior mistress," he enjoined her. "Tell her that
        Ch'ing Wen got a slight chill yesterday. That as it's nothing to speak of,
        and Hsi Jen is besides away, there would be, more than ever, no one here to
        look after things, were she to go home and attend to herself, so let her
        send for a doctor to come quietly by the back entrance and see what's the
        matter with her;}
}
{
}

\Verse{56}
{
    \zA{}
}
{
    \eA{but don't let her breathe a word about it to Madame Wang, my mother." The
        old nurse was away a considerable time on the errand. On her return, "Our
        senior mistress," she reported, "has been told everything. She says that:
        'if she gets all right, after taking a couple of doses of medicine, it will
        be well and good. But that in the event of not recovering, it would, really,
        be the right thing for her to go to her own home. That the season isn't
        healthy at present, and that if the other girls caught her complaint it
        would be a small thing; but that the good health of the young ladies is a
        vital matter.'"}
}
{
}

\Verse{57}
{
    \zA{}
}
{
    \eA{Ch'ing Wen was lying in the winter apartment, coughing and coughing, when
        overhearing (Li Wan's) answer, she lost control over her temper. "Have I got
        such a dreadful epidemic," she said, "that she fears that I shall bring it
        upon others? I'll clear off at once from this place; for mind you don't get
        any headaches and hot heads during the course of your lives." "While
        uttering her grievances, she was bent upon getting up immediately, when
        Pao-yü hastened to smile and to press her down. "Don't lose your temper," he
        advised her. "This is a responsibility which falls upon her shoulders, so
        she is afraid lest Madame Wang might come to hear of it, and call her to
        task. She only made a harmless remark.}
}
{
}

\Verse{58}
{
    \zA{}
}
{
    \eA{But you've always been prone to anger, and now, as a matter of course your
        spleen is larger than ever." But in the middle of his advice to her, a
        servant came and told him that the doctor had arrived. Pao-yü accordingly
        crossed over to the off side, and retired behind the bookcase; from whence
        he perceived two or three matrons, whose duty it was to keep watch at the
        back door, usher the doctor in. The waiting-maids, meanwhile, withdrew out
        of the way. Three or four old nurses dropped the deep-red embroidered
        curtain, suspended in the winter apartment. Ch'ing Wen then simply stretched
        out her hand from among the folds of the curtain.}
}
{
}

\Verse{59}
{
    \zA{}
}
{
    \eA{But the doctor noticed that on two of the fingers of her hand, the nails,
        which measured fully two or three inches in length, still bore marks of the
        pure red dye from the China balsam, and forthwith he turned his head away.
        An old nurse speedily fetched a towel and wiped them for her, when the
        doctor set to work and felt her pulse for a while, after which he rose and
        walked into the outer chamber. "Your young lady's illness," he said to the
        old nurses, "arises from external sources, and internal obstructive
        influences, caused by the unhealthiness of the season of late. Yet it's only
        a slight chill, after all. Fortunately, the young lady has ever been
        moderate in her drinking and eating. The cold she has is nothing much. It's
        mainly because she has a weak constitution that she has unawares got a bit
        of a chill.}
}
{
}

\Verse{60}
{
    \zA{}
}
{
    \eA{But if she takes a couple of doses of medicine to dispel it with, she'll be
        quite right." So saying, he followed once more the matron out of the house.
        Li Wan had, by this time, sent word to the various female domestics at the
        back entrance, as well as to the young maids in the different parts of the
        establishment to keep in retirement. All therefore that the doctor perceived
        as he went along was the scenery in the garden. But not a single girl did he
        see. Shortly, he made his exit out of the garden gate, and taking a seat in
        the duty-lodge of the servant-lads, who looked after the garden-entrance, he
        wrote a prescription. "Sir," urged an old nurse, "don't go yet. Our young
        master is fretful and there may be, I fancy, something more to ask you."}
}
{
}

\Verse{61}
{
    \zA{}
}
{
    \eA{"Wasn't the one I saw just now a young lady," the doctor exclaimed with
        eagerness, "but a young man, eh? Yet the rooms were such as are occupied by
        ladies. The curtains were besides let down. So how could the patient I saw
        have ever been a young man?" "My dear sir," laughed the old nurse, "it isn't
        strange that a servant-girl said just now that a new doctor had been sent
        for on this occasion, for you really know nothing about our family matters.
        That room is that of our young master, and that is a girl attached to the
        apartments; but she's really a servant-maid. How ever were those a young
        lady's rooms?}
}
{
}

\Verse{62}
{
    \zA{}
}
{
    \eA{Had a young lady fallen ill, would you ever have penetrated inside with such
        ease?" With these words, she took the prescription and wended her way into
        the garden. When Pao-yü came to peruse it, he found, above, such medicines
        mentioned as sweet basil, platycodon, carraway seeds, mosla dianthera, and
        the like; and, below, citrus fusca and sida as well. "He deserves to be
        hanged! He deserves death!" Pao-yü shouted. "Here he treats girls in the
        very same way as he would us men! How could this ever do? No matter what
        internal obstruction there may be, how could she ever stand citrus and sida?
        Who asked him to come? Bundle him off at once; and send for another, who
        knows what he's about." "Whether he uses the right medicines or not," the
        old nurse pleaded, "we are not in a position to know.}
}
{
}

\Verse{63}
{
    \zA{}
}
{
    \eA{But we'll now tell a servant-lad to go and ask Dr. Wang round. It's easy
        enough! The only thing is that as this doctor wasn't sent for through the
        head manager's office his fee must be paid to him." "How much must one give
        him?" Pao-yü inquired. "Were one to give him too little, it wouldn't look
        nice," a matron ventured. "He should be given a tael. This would be quite
        the thing with such a household as ours." "When Dr. Wang comes," Pao-yü
        asked, "how much is he given?" "Whenever Dr. Wang and Dr. Chang come," a
        matron smilingly explained, "no money is ever given them. At the four
        seasons of each year however presents are simply sent to them in a lump.
        This is a fixed annual custom. But this new doctor has come only this once
        so he should be given a tael."}
}
{
}

\Verse{64}
{
    \zA{}
}
{
    \eA{After this explanation, Pao-yü readily bade She Yüeh go and fetch the money.
        "I can't make out where sister Hua put it;" She Yüeh rejoined. "I've often
        seen her take money out of that lacquered press, ornamented with designs
        made with shells;" Pao-yü added; "so come along with me, and let's go and
        search." As he spoke, he and She Yüeh came together into what was used as a
        store-room by Hsi Jen. Upon opening the shell-covered press, they found the
        top shelf full of pens, pieces of ink, fans, scented cakes, various kinds of
        purses, handkerchiefs and other like articles, while on the lower shelf were
        piled several strings of cash. But, presently they pulled out the drawer,
        when they saw, in a small wicker basket, several pieces of silver, and a
        steelyard. She Yüeh quickly snatched a piece of silver.}
}
{
}

\Verse{65}
{
    \zA{}
}
{
    \eA{Then raising the steelyard, "Which is the one tael mark?" she asked. Pao-yü
        laughed. "It's amusing that you should appeal to me!" he said. "You really
        behave as if you had only just come!" She Yüeh also laughed, and was about
        to go and make inquiries of some one else, when Pao-yü interfered. "Choose a
        piece out of those big ones and give it to him, and have done," he said. "We
        don't go in for buying and selling, so what's the use of minding such
        trifles!" She Yüeh, upon hearing this, dropped the steelyard, and selected a
        piece, which she weighed in her hand. "This piece," she smiled, "must, I
        fancy, be a tael. But it would be better to let him have a little more.
        Don't let's give too little as those poor brats will have a laugh at our
        expense.}
}
{
}

\Verse{66}
{
    \zA{}
}
{
    \eA{They won't say that we know nothing about the steelyard; but that we are
        designedly mean." A matron who stood at the threshold of the door, smilingly
        chimed in. "This ingot," she said, "weighs five taels. Even if you cut half
        of it off, it will weigh a couple of taels, at least. But there are no sycee
        shears at hand, so, miss, put this piece aside and choose a smaller one."
        She Yüeh had already closed the press and walked out. "Who'll go and fumble
        about again?" she laughed. "If there's a little more, well, you take it and
        finish." "Be quick," Pao-yü remarked, "and tell Pei Ming to go for another
        doctor. It will be all right." The matron received the money and marched off
        to go and settle matters. Presently, Dr. Wang actually arrived, at the
        invitation of Pei Ming.}
}
{
}

\Verse{67}
{
    \zA{}
}
{
    \eA{First and foremost he felt the pulse and then gave the same diagnosis of the
        complaint (as the other doctor did) in the first instance. The only
        difference being that there was, in fact, no citrus or sida or other similar
        drugs, included in the prescription. It contained, however, false
        sarsaparilla roots, dried orange peel, peonia albifora, and other similar
        medicines. But the quantities were, on the other hand, considerably smaller,
        as compared with those of the drugs mentioned in the former prescription.
        "These are the medicines," Pao-yü ejaculated exultingly, "suitable for
        girls! They should, it's true, be of a laxative nature, but never over and
        above what's needful.}
}
{
}

\Verse{68}
{
    \zA{}
}
{
    \eA{When I fell ill last year, I suffered from a chill, but I got such an
        obstruction in the viscera that I could neither take anything liquid or
        substantial, yet though he saw the state I was in, he said that I couldn't
        stand sida, ground gypsum, citrus and other such violent drugs. You and I
        resemble the newly-opened white begonia, Yün Erh sent me in autumn. And how
        could you resist medicines which are too much for me? We're like the lofty
        aspen trees, which grow in people's burial grounds. To look at, the branches
        and leaves are of luxuriant growth, but they are hollow at the core." "Do
        only aspen trees grow in waste burial grounds?" She Yüeh smiled. "Is it
        likely, pray, that there are no fir and cypress trees? What's more loathsome
        than any other is the aspen. For though a lofty tree, it only has a few
        leaves;}
}
{
}

\Verse{69}
{
    \zA{}
}
{
    \eA{and it makes quite a confused noise with the slightest puff of wind! If you
        therefore deliberately compare yourself to it, you'll also be ranging
        yourself too much among the common herd!" "I daren't liken myself to fir or
        cypress;" Pao-yü laughingly retorted. "Even Confucius says: 'after the
        season waxes cold, one finds that the fir and cypress are the last to lose
        their foliage,' which makes it evident that these two things are of high
        excellence. Thus it's those only, who are devoid of every sense of shame,
        who foolishly liken themselves to trees of the kind!"}
}
{
}

\Verse{70}
{
    \zA{}
}
{
    \eA{While engaged in this colloquy, they perceived the old matron bring the
        drugs, so Pao-yü bade her fetch the silver pot, used for boiling medicines
        in, and then he directed her to prepare the decoction on the brasier. "The
        right thing would be," Ch'ing Wen suggested, "that you should let them go
        and get it ready in the tea-room; for will it ever do to fill this room with
        the smell of medicines?" "The smell of medicines," Pao-yü rejoined, "is far
        nicer than that emitted by the whole lot of flowers. Fairies pick medicines
        and prepare medicines. Besides this, eminent men and cultured scholars
        gather medicines and concoct medicines; so that it constitutes a most
        excellent thing. I was just thinking that there's everything and anything in
        these rooms and that the only thing that we lack is the smell of medicines;}
}
{
}

\Verse{71}
{
    \zA{}
}
{
    \eA{but as luck would have it, everything is now complete." Speaking, he lost no
        time in giving orders to a servant to put the medicines on the fire. Next,
        he advised She Yüeh to get ready a few presents and bid a nurse take them
        and go and look up Hsi Jen, and exhort her not to give way to excessive
        grief. And when he had settled everything that had to be seen to, he
        repaired to the front to dowager lady Chia's and Madame Wang's quarters, and
        paid his respects and had his meal. Lady Feng, as it happened, was just
        engaged in consulting with old lady Chia and Madame Wang.}
}
{
}

\Verse{72}
{
    \zA{}
}
{
    \eA{"The days are now short as well as cold," she argued, "so wouldn't it be
        advisable that my senior sister-in-law, Mrs. Chia Chu, should henceforward
        have her repasts in the garden, along with the young ladies? When the
        weather gets milder, it won't at all matter, if they have to run backward
        and forward." "This is really a capital idea!" Madame Wang smiled. "It will
        be so convenient during windy and rainy weather. To inhale the chilly air
        after eating isn't good. And to come quite empty, and begin piling up a lot
        of things in a stomach full of cold air isn't quite safe.}
}
{
}

\Verse{73}
{
    \zA{}
}
{
    \eA{It would be as well therefore to select two cooks from among the women, who
        have, anyhow, to keep night duty in the large five-roomed house, inside the
        garden back entrance, and station them there for the special purpose of
        preparing the necessary viands for the girls. Fresh vegetables are subject
        to some rule of distribution, so they can be issued to them from the general
        manager's office. Or they might possibly require money or be in need of some
        things or other. And it will be all right if a few of those pheasants, deer,
        and every kind of game, be apportioned to them." "I too was just thinking
        about this," dowager lady Chia observed. "The only thing I feared was that
        with the extra work that would again be thrown upon the cook-house, they
        mightn't have too much to do."}
}
{
}

\Verse{74}
{
    \zA{}
}
{
    \eA{"There'll be nothing much to do," lady Feng replied. "The same apportionment
        will continue as ever. In here, something may be added; but in there
        something will be reduced. Should it even involve a little trouble, it will
        be a small matter. If the girls were exposed to the cold wind, every one
        else might stand it with impunity; but how could cousin Lin, first and
        foremost above all others, resist anything of the kind? In fact, brother Pao
        himself wouldn't be proof against it. What's more, none of the various young
        ladies can boast of a strong constitution." What rejoinder old lady Chia
        made to lady Feng, at the close of her representations, is not yet
        ascertained;}
}
{
}

\Verse{75}
{
    \zA{}
}
{
    \eA{so, reader, listen to the explanations you will find given in the next
        chapter.}
}
{
}
